\begin{chineseabstract}
存内计算（Compute-in-Memory, CIM）架构通过在存储阵列内部直接执行矩阵--向量乘法，消除了制约传统数字系统的数据搬运瓶颈。
在诸多新兴器件路线中，有机光电子器件凭借其多级电导可调性、低压操作与本征光敏性，为神经形态视觉系统提供了独特的物理载体。
然而，将器件级进展转化为系统级视觉性能，仍面临器件变异性、有限模数转换器（Analog-to-Digital Converter, ADC）精度、保留漂移及跨制造实例可迁移性等关键挑战。

本研究建立了一个面向有机光电子存内计算阵列的可复用硬件感知训练（Hardware-Aware Training, HAT）仿真框架，
支持具有真实器件非理想性的视觉 Transformer（Vision Transformer, ViT）与卷积神经网络（Convolutional Neural Network, CNN）。
该框架采用配置文件驱动设计，可将文献报道的器件特性（电导窗口、周期到周期变异性、器件到器件失配、保留动力学、光响应非线性、写入非线性）映射为任务级视觉精度。

实验首先系统评估了 ResNet-18、ConvNeXt-Tiny 与 Tiny-ViT-5M 在 CIFAR-10、CIFAR-100、Flowers-102 及 ImageNet-1k 上的 baseline 精度。
在此基础上，本研究提出 HAT 训练范式的系统分类学，并发现\emph{硬件实例过拟合}现象：
标准固定掩膜 HAT 在训练时硬件实例上表现良好，但在全新器件失配实例上崩溃至接近随机水平（CIFAR-10 上 $10.00\pm0.00\%$）。
为应对此问题，本研究提出 Ensemble HAT——在每个训练轮次边界重采样结构化器件到器件（Device-to-Device, D2D）失配场——
将 fresh-instance 精度恢复至 $86.16\pm0.19\%$。

进一步地，本研究构建了面向有机光电子 CIM 部署的\emph{失效模式图谱}（failure mode atlas），
系统刻画了五种可复现的失效模式及其边界：
（1）硬件实例过拟合与 Ensemble HAT 缓解；
（2）严重写入非线性（$NL=2.0$）下的双注意力通路崩溃；
（3）多层感知机（Multi-Layer Perceptron, MLP）线性补偿的 $32\%$ fresh-instance 诊断上界；
（4）空间相关 D2D 的有界退化；
（5）保留漂移的 $79\%$ 渐近平台期。
所有失效模式均通过严格的负结果（negative results）加以界定，从而将瓶颈定位为结构性泛化障碍而非优化缺口。

本研究的贡献包括：
（1）一个可复用、可扩展的有机光电子 CIM 仿真框架；
（2）HAT 范式的系统分类学与硬件实例过拟合的发现；
（3）一个受经验数据约束的部署包络，明确了哪些噪声模型与训练干预组合能够迁移至新硬件实例、哪些不能。
这些结果为有机光电子神经形态硬件的算法--器件协同设计提供了系统级决策依据。
\chinesekeywordstype{存内计算；硬件感知训练；视觉 Transformer；有机光电子器件；硬件实例过拟合}{应用基础研究}
\end{chineseabstract}
